\section{Introduction}\label{sec:intro}\mymarginpar{Meta 1, R1 O2, R3 O1}

Succinct data structures have become essential in applications like large-scale
databases, graph processing, and text indexing, especially in memory-constrained
scenarios~\cite{sds_application}. They are string-data representations that use
close to the theoretical minimum number of bits while still supporting efficient
queries~\cite{sds_overview, munro2018succinct}. We focus on \defn{succinct
  tries}~\cite{surf, pdt, coco_0, coco_1, marisa}, which match or exceed the
performance of regular tries \rev{with orders of magnitude less memory.}

Succinct tries separate the trie \defn{topology} from the string \defn{data},
encoding the topology in a compact bitvector that uses only about 2 bits per
node on average compared to at least 128 bits for two pointers in traditional
tries~\cite{jacobson1989space}. Navigation requires auxiliary index operations
(e.g., ``\rankop'' and ``\selectop'') on top of this bitvector.

\begin{table}[t]
  \caption{ Average number of LLC misses per query on the two largest datasets.
    FST~\cite{surf}, CoCo-trie~\cite{coco_1} and Marisa~\cite{marisa} are
    \bv-based succinct tries, while C-ART~\cite{hybrid_index} is a pointer-based
    compact trie.  Tries prefixed with \csquared- are our cache-optimized
    versions.  }
  \resizebox{\columnwidth}{!}{%
    \begin{tabular}{llllllll}
        \toprule
        \textit{Dataset}
        & \textit{FST}
        & \textit{\begin{tabular}{@{}l@{}}\csquared-\\FST\end{tabular}}
        & \textit{CoCo'}
        & \textit{\begin{tabular}{@{}l@{}}\csquared-\\CoCo\end{tabular}}
        & \textit{Marisa}
        & \textit{\begin{tabular}{@{}l@{}}\csquared-\\Marisa\end{tabular}}
        & \textit{C-ART} \\
        \midrule

        \wikidata~\cite{wiki-dataset}
        & 82
        & 78
        & 66
        & 49
        & 56
        & 58
        & \textbf{23} \\

        \logdata~\cite{log-dataset}
        & 235
        & 230
        & 57
        & \textbf{45}
        & 58
        & 62
        & 117 \\

        \bottomrule
    \end{tabular}
}
\label{tab:cache-motivation}
\end{table}
%%% Local Variables:
%%% mode: latex
%%% TeX-master: "../main"
%%% End:


\newmymarginpar{Meta 1, R1 O2, R3 O1/D1}\para{Challenges to locality}
  Existing \bv-based succinct tries exhibit suboptimal locality in two respects:
  1) the \emph{bitvector layout} in the topology, and 2) the \newmymarginpar{R2
    O3/D6}\newrev{\defn{unary paths}, or maximal trie paths of at least 2 nodes
    in which every node except the last has exactly one child}, in the data.
  For example, in~\figref{louds-example}, paths 0-1-4 (``ca'') and 0-2-5
  (``su'') are internal unary paths, whereas path 7-14-18 (``ie'') is a suffix
  unary path.~\secreftwo{bv-design}{unary-path} detail both challenges.

\rev{\tabref{cache-motivation}} confirms locality challenges due to navigation
operations across state-of-the-art succinct tries. For example, a child
navigation in Marisa~\cite{marisa} incurs at least 3 cache misses to access the
topology \bv compared to 1 in a pointer-based trie to follow the pointer.

\begin{table}[t]

\caption{
Unary-path statistics for the two largest datasets.
\#Branch denotes the number of branching edges. Given a path length $\ell$, we
report the percentage of paths with lengths
 $\ell=1$, $1<\ell \leq 3$, and  %\% $1<\ell<\ell_{\mathrm{cut}} = 3$), and
$\ell> 3$.
% correspond to the three unary-path categories defined in Section~\ref{sec:unary}.
$\ell_{AVG}$ and $\ell_{MAX}$ denote the average and maximum
lengths of compressible unary paths. %(links).
}
\centering
\small
\begin{tabular}{lrrrrrr}
\toprule
\textit{Dataset}
& \textit{\# Branch}
& \textit{$\ell = 1$}
& \textit{$1 < \ell \leq 3$}
& \textit{$\ell > 3$}
& $\ell_{AVG}$
& $\ell_{MAX}$ \\
\midrule

\wikidata~\cite{wiki-dataset}
& 467K
& 15.5\%
& 0.4\%
& 84.1\%
& 271
& 8676 \\

\logdata~\cite{log-dataset}
& 7,284K
& 17.4\%
& 3.8\%
& 78.8\%
& 65
& 8196 \\

\bottomrule
\end{tabular}

\label{tab:unary-wiki-log}

\end{table}

%%% Local Variables:
%%% mode: latex
%%% TeX-master: "../main"
%%% End:


\newnew{Additionally, practical datasets often include strings with long
  shared prefixes followed by diverse \new{dangling} suffixes, which naturally
  introduce large numbers of unary paths in trie-based representations.  For
  example, collections of hyperlinks frequently share common prefixes such as
  ``en.wikipedia.org/wiki/'', while differing only in their trailing titles. As
  a result, large portions of the trie structure degenerate into long chains of
  unary nodes.
  %

  Without effective compression, unary paths can consume a dominant fraction of
  succinct-trie storage, increasing memory footprint and disrupting cache
  locality.  \newmymarginpar{R1 O1}For example, as shown
  in~\tabref{unary-wiki-log}, the majority of suffix regions in realistic
  datasets consist of long, redundant unary chains in the two largest tested
  inputs. In the \wikidata\ dataset, 84.1\% of all branch edges correspond to
  compressible unary paths, with an average compressible length of 271
  characters and a maximum compressible length of over 8K characters. Similarly,
  the \logdata\ dataset contains 78.8\% compressible unary paths.}


% \para{Cache-conscious \bv redesign}
\para{Optimizing for locality in \bv-based succinct tries} \emph{We introduce
  \csquared, a set of techniques that improve locality in both the topology and
  the data of succinct tries.} \newmymarginpar{Meta 3, R1 O3, R3 O2}\newrev{We target two operations: the \defn{existence query} (checking
  whether a key exists) and the \defn{range query} (returning a range of keys).}

The first ``C'' of~\csquared improves cache locality of the \selectop index
in succinct tries. We introduce the \defn{functional index}, a novel
\selectop index layout that maps straightforwardly to bit positions, enabling
both \rankop and \selectop indexes to be inlined in the bit sequence.
%
\mymarginpar{Meta 4}\rev{Using the functional index, we cut random accesses to
  two per child navigation in LOUDS-based succinct tries, which include many
  state-of-the-art tries such as the FST~\cite{surf},
  CoCo-trie~\cite{coco_0}, and Marisa~\cite{marisa}. \tabref{cache-motivation}
  and \secref{experiments} show that this optimization reduces cache misses by
  19\% on average with at most $2\%$ overall space overhead.}

\mymarginpar{R3 O2, Meta 4}\newrev{The second ``C'' of~\csquared is an
  \defn{adaptive compression} algorithm that gives \emph{any} succinct trie
  access to compression in the ``tail container,'' which stores the remaining
  string suffixes at the end of trie paths. Building on
  ``\repair''~\cite{larsson2000off} and ``recursion''~\cite{marisa}, two
  orthogonal techniques for path compression, we enable all
  tries to use both recursion and the Fast Static Symbol Table
  (\fsst)~\cite{fsst}, a lightweight scheme that achieves good compression
  ratios with fast build times. Although recursion and compressed tail
  containers appear in prior work, this is the first generalization of
  recursion across tail containers, exposing space-time tradeoffs for any
  succinct trie.}

% \para{Applying \csquared to state-of-the-art succinct tries}



\begin{figure}[t]
    \pgfplotsset{every axis title/.append style={at={(0.5,0.95)}}}
    \centering
    \begin{tikzpicture}
    \begin{groupplot}[
        % group style={group size=2 by 1, horizontal sep=35pt},
        group style={group size=1 by 2},
        width=8cm, height=4cm,
        ymajorgrids=true,
        xmajorgrids=true,
        minor x tick num=1,
        xminorgrids=true,
        minor y tick num= 3,
        yminorgrids = true,
        minor grid style=loosely dotted,
        legend columns=4,
        legend entries={FST, \bvdesign-FST, \csquared-FST, CoCo', \bvdesign-CoCo, \csquared-CoCo, Marisa, \bvdesign-Marisa, \csquared-Marisa, Marisa-1, \bvdesign-Marisa-1, \csquared-Marisa-1, PDT, C-ART},
        legend to name=PerformanceLegend,
        scatter/classes={
            FST={mark=*, fill=safe-brown},
            C1-FST={mark=*, fill=safe-olive},
            C2-FST={mark=*, fill=safe-teal},
            CoCo'={mark=rotated triangle*, safe-black},
            C1-CoCo={mark=rotated triangle*, safe-lavender},
            C2-CoCo={mark=rotated triangle*, safe-peach},
            %C2-CoCo={mark=square*, orange,rotate=45, mark size=3pt},
            Marisa={mark=diamond*, safe-plum},
            C1-Marisa={mark=diamond*, safe-green},
            C2-Marisa={mark=diamond*, safe-sky},
            Marisa-1={mark=halfsquare*, safe-cerulean},
            C1-Marisa-1={mark=halfsquare*, safe-brick},
            C2-Marisa-1={mark=halfsquare*, orange},
            % C2-FST'={mark=halfsquare*, safe-teal},
            % C2-CoCo'={mark=triangle*, red},
            % C2-Mar'={mark=diamond*, blue},
            % C2-FST''={mark=halfsquare*, lightgray},
            % C2-CoCo''={mark=triangle*, lightgray},
            % C2-Mar''={mark=diamond*, lightgray},
            % FST={mark=x, black},
            % CoCo={mark=square*, violet},
            % Marisa={mark=+, magenta},
            % Mar-1={mark=otimes*, orange},
            % Mar-2={mark=triangle, violet},
            % PDT={mark=Mercedes star flipped, violet},
            PDT={mark=otimes, blue},
            C-ART={mark=oplus, red}
        },
        scatter src=explicit symbolic
    ]

    \nextgroupplot[
        title={WIKI (359 MB)},
        % xlabel={latency (\textit{ns})},
        ylabel={Size (normalized to ART)},
        legend to name=PerformanceLegend
    ]
    \addplot[scatter,only marks,mark size=3.5pt,point meta=explicit,scatter src=explicit symbolic]
    table[x=query,y=space,meta=trie]{tsvs/results-wiki-new-2.tsv};

    \nextgroupplot[
        title={LOG (586 MB)},
        xlabel={Query latency (normalized to ART)},
        ylabel={Size (normalized to ART)},
        legend to name=PerformanceLegend
    ]
    \addplot[scatter,only marks,mark size=3.5pt,point meta=explicit,scatter src=explicit symbolic]
    table[x=query,y=space,meta=trie]{tsvs/results-log-new-2.tsv};

    \end{groupplot}
    \end{tikzpicture}

    \begin{center}
    \ref{PerformanceLegend}
    \end{center}
    \caption{
    % \helen{Fix legend to fit in one column}
    \rev{Space-time performance
        comparison on \wikidata~\cite{wiki-dataset} and
        \logdata~\cite{log-dataset} (normalized to ART). We use FSST as the tail
        container in all \csquared-tries. The prefixes "\bvdesign-"/"\csquared-"
        indicate different optimization strategies (bitvector redesign /
        bitvector redesign + unary path compression); the suffix "-1" means we
        enforce one recursion (see \secref{experiments}). Lower and to the left
        is better.}}
    \label{fig:intro-results-summary}
\end{figure}
%%% Local Variables:
%%% mode: latex
%%% TeX-master: "../sample-sigconf"
%%% End:

\para{Contributions} Our primary contributions are as follows:
\begin{itemize}[leftmargin=*,noitemsep]
\item We redesign the succinct-trie \rev{topology} with a \emph{cache-conscious}
  functional index that inlines both \rankop and \selectop indexes, minimizing
  cache misses during trie navigation.
\item \rev{We propose an \emph{adaptive compression scheme} that selects the
    best compression strategy per dataset to balance space and query
    performance.} \newrev{Our experiments indicate that the chosen strategy improves
  space efficiency by \toplevelspacesavings on average with similar
   (within $1.05\times$) query performance compared to the original succinct trie versions.} Furthermore, increasing compression aggressiveness exposes further space-time tradeoffs.
  \newrev{For example, on Marisa, increasing ``recursion'' for more compression yields
  \marisarecursionspace space savings, but \marisarecursiontime slower queries.}
\item We integrate these techniques into three state-of-the-art succinct tries
  --- FST~\cite{surf}, CoCo-trie~\cite{coco_0, coco_1}, and Marisa~\cite{marisa}
  ---
  %reducing cache misses \rev{during}\mymarginpar{Meta 5, R2 D17} query
  %operations by \rev{\todo{update ratio} on average} and
  improving query
  performance on the whole (with both \bvdesign and \compressionscheme
  optimizations) by \csqfstspeedup, \csqcocospeedup, and \csqmarisaspeedup,
  respectively.
\end{itemize}


\para{Results summary} Overall, \csquared-optimized
tries improve cache locality and performance over their original counterparts with negligible (<4\%) additional space overhead. 
\csquared-FST achieves better space-time performance than FST on all datasets. 
\rev{ 
    % \csquared-CoCo outperforms CoCo' by \csqcocospeedup in query time and \csqcocospacesaving in space
    \csquared-CoCo outperforms CoCo' by \csqcocospeedup in query latency and is \csqcocospacesaving smaller in space, on  average across the six datasets
        \footnote{The original CoCo-trie was designed for prefix-only datasets and fails to build on the original datasets. 
        We implemented CoCo' by integrating CoCo-trie's topology with \csquared-CoCo.~\secref{experiments} provides the full details and shows that CoCo and CoCo' achieve almost identical performance and space usage on the prefix-only datasets.
        }. 
    \csquared-Marisa improves Marisa's query performance by \csqmarisaspeedup with similar memory consumption, and is generally the strongest \csquared index.
}

\figref{intro-results-summary} \rev{shows the query latency and size of the
  different tries} on the \rev{\wikidata} and \rev{\logdata} datasets.
\csquared-Marisa dominates all other succinct tries on query latency with the
lowest (or within a few percent of the lowest) space usage. \rev{ART and
C-ART, two non-succinct tries, achieve lower query latency at the cost of
higher space usage.}

\rev{These space-time improvements stem from the cache-locality gains of
  \csquared, as demonstrated by the cache-miss counts
  in~\tabref{cache-motivation}.}

%%% Local Variables:
%%% mode: latex
%%% TeX-master: "main"
%%% End:
